\documentclass[11pt,a4paper]{article}
\usepackage[margin=26mm]{geometry}
\usepackage[T1]{fontenc}
\usepackage{lmodern,microtype}
\usepackage{amsmath,amssymb,amsthm,mathtools,booktabs}
\usepackage[hidelinks]{hyperref}
\hypersetup{pdftitle={Lossless Assembly of an Eight-Point Certificate},
 pdfauthor={RH-Weil research project}}
\newtheorem{theorem}{Theorem}[section]
\newtheorem{lemma}[theorem]{Lemma}
\newtheorem{proposition}[theorem]{Proposition}
\newtheorem{inputclaim}[theorem]{Computational input}
\theoremstyle{remark}
\newtheorem{remark}[theorem]{Remark}
\DeclareMathOperator{\tr}{tr}
\DeclareMathOperator{\sinc}{sinc}
\newcommand{\R}{\mathbb R}
\newcommand{\I}{[-1/2,1/2]}
\title{Lossless Assembly of an Eight-Point Certificate\\
for Simple Critical Zeros of the Riemann Zeta Function}
\author{RH--Weil research project}
\date{October 8, 2026}
\begin{document}
\maketitle
\begin{abstract}
We combine an existing seven-gap certificate for a perturbed
Montgomery--Taylor window with the close-pair separation mechanism used
by Knausg{\aa}rd. An explicit band-limited majorant has mass strictly
below two and dominates the same perturbed window. It removes the
spectral truncation loss on the separated set, while disjoint close pairs
pay the same per-point reward. Sliding the existing eight-point
certificate directly over the separated set gives
\[
 \liminf_{T\to\infty}\frac{N_0^s(T)}{N(T)}
 \ge \frac{66812491}{99194997}
 =0.6735469834229644\ldots .
\]
Here the denominator counts all nontrivial zeros with multiplicity,
and the numerator counts simple zeros on the critical line.
The analytic argument is unconditional and retains off-line zeros.
The computational input is the previously admitted PC8CL continuous
seven-gap certificate. We give a separate, small exact rational
certificate for the new majorant and the energy constant. This paper
does not claim a complete Lean kernel verification or external
peer-review certification.
\end{abstract}

\section{Statement, provenance and comparison}

Let $N(T)$ count the nontrivial zeros $\rho$ of $\zeta$, with
$0<\Im\rho\le T$, including multiplicity. Let $N_0^s(T)$ count the
simple zeros among them with $\Re\rho=1/2$.

\begin{theorem}\label{thm:main}
Using the continuous PC8CL certificate stated in
Input~\ref{in:pc8}, one has, unconditionally,
\begin{equation}\label{eq:main}
 \liminf_{T\to\infty}\frac{N_0^s(T)}{N(T)}
 \ge \frac{66812491}{99194997}.
\end{equation}
No zero-free strip, Riemann hypothesis or higher correlation
conjecture is assumed.
\end{theorem}

The local certificate and the window below are existing inputs from
the fixed public AM eight-point construction \cite{PC8}.
Their prior square-root spectral assembly gave the admitted rational
bound $118513839290/175971686899$. The present change is in the
assembly: we use a band-limited majorant, remove disjoint close pairs,
and apply the old local inequality without a square-root loss.
We adapt the simple-point separation mechanism in Knausg{\aa}rd
\cite[Section~5]{Knaus}, whose v1 uses a different window and local
certificate to obtain $1669159/2478195$.
We reprove the finite and analytic transfer needed here.
Our exact comparison with that stated v1 benchmark is
\begin{equation}\label{eq:comparison}
 \frac{66812491}{99194997}-\frac{1669159}{2478195}
 =\frac{719712074}{81941515196805}>0.
\end{equation}
The increase is approximately $0.0008783241$ percentage points.
This comparison is with a checked, fixed benchmark; it does not
establish exhaustive priority over all unpublished or concurrent work.
The proportion in \eqref{eq:main} is a zero-counting lower bound,
rather than a measure of progress toward a proof of RH.

\section{Window and admitted local certificate}

Put $\beta=1/\sqrt2$ and $Z_0=\sqrt2\sin\beta$.
On $I=\I$ define
\begin{equation}\label{eq:window}
 f(u)=\frac{\cos(\sqrt2u)+\sum_{j=1}^{12}c_j\cos(2\pi ju)}{Z_0},
 \qquad K(t)=\int_I f(u)e^{-2\pi itu}\,du.
\end{equation}
Extend $f$ by zero outside $I$. The integer numerators $10^9c_j$ are
\[
\begin{split}
(&12310798,-15041681,3867664,6489926,-992327,5580097,\\
 &-6846472,3781297,-5670353,3355089,-450523,-218483).
\end{split}
\]
Thus $S=\sum_j|c_j|=6460471/10^8<13/200$.
Since $\cos(\sqrt2u)\ge1-u^2\ge3/4$ on $I$, the numerator
is positive. The integer cosine modes integrate to zero, so
$f$ is an even probability density. Also
\begin{equation}\label{eq:normalization}
 Z_0=\int_I\cos(\sqrt2u)\,du\ge 11/12.
\end{equation}

For eight ordered points $y_0,\ldots,y_7$, write
$g_r=y_{r+1}-y_r\ge0$, $0\le r\le6$. Set
\[
 c=\frac{805003}{10^8},\qquad B=\sum_rb_r=\frac{404350}{10^8},
\]
\[
 (10^8b_r)_{r=0}^6=(28898,57272,75526,80958,75526,57272,28898).
\]
All pair weights $a_{ij}$ are nonnegative. Table~\ref{tab:weights}
gives their integer numerators.

\begin{table}[ht]
\centering\small
\begin{tabular}{c*{7}{r}}
\toprule
$i\backslash j$&1&2&3&4&5&6&7\\
\midrule
0&13630830&0&19565904&20317441&45030671&100000000&200000000\\
1&&29997996&30621878&47766489&79682558&109938656&100000000\\
2&&&37356140&69378121&65335210&79682558&45030671\\
3&&&&38030065&69378121&47766489&20317441\\
4&&&&&37356140&30621878&19565904\\
5&&&&&&29997996&0\\
6&&&&&&&13630830\\
\bottomrule
\end{tabular}
\caption{The weights $10^8a_{ij}$ of the fixed AM certificate.}
\label{tab:weights}
\end{table}

\begin{inputclaim}[Fixed PC8CL certificate]\label{in:pc8}
For every $(g_0,\ldots,g_6)\in[0,\infty)^7$,
\begin{equation}\label{eq:local}
 \sum_{r=0}^6b_rg_r+
 \sum_{0\le i<j\le7}a_{ij}
 K\left(\sum_{r=i}^{j-1}g_r\right)^2\ge c.
\end{equation}
\end{inputclaim}

This is a continuous real-gap assertion, not a grid sampling claim.
Its provenance, finite replay and continuous soundness evidence are
specified in Section~\ref{sec:trust}. We use the complete admitted
certificate, rather than the upstream final zeta theorem.
Direct integer summation of Table~\ref{tab:weights} gives
\begin{equation}\label{eq:capacity}
\begin{split}
10^8\left(\sum_{i=0}^{7-s}a_{i,i+s}\right)_{s=1}^7
={}&(199999997,199999998,199999996,199999998,\\
 &199999998,200000000,200000000).
\end{split}
\end{equation}
In particular each span-length budget is at most two.

\begin{lemma}[Direct sliding]\label{lem:sliding}
For any $m$ ordered real points $y_0\le\cdots\le y_{m-1}$,
with span taken to be zero when $m\le1$,
\begin{equation}\label{eq:sliding}
 \sum_{i\ne j}K(y_i-y_j)^2
 \ge c(m-7)-B\,\operatorname{span}(y).
\end{equation}
\end{lemma}
\begin{proof}
For $m\ge8$, sum \eqref{eq:local} over the $m-7$ consecutive
eight-point windows. A fixed pair of index distance $s\le7$ receives
total weight at most $\sum_i a_{i,i+s}\le2$; pairs farther apart
receive no weight. A fixed adjacent gap receives total coefficient
at most $\sum_rb_r=B$. All terms are nonnegative, so the weighted
pair sum is at most the full ordered pair sum in
\eqref{eq:sliding}. This proves the assertion. For $m\le7$ its
right side is nonpositive.
\end{proof}

\section{A band-limited majorant for the same window}

We use $\sinc t=\sin(\pi t)/(\pi t)$, with $\sinc0=1$, and put
\begin{equation}\label{eq:majorant}
 \theta=\frac45,\quad \gamma=\frac{61}{100},\qquad
 g(u)=\gamma\{\sinc(\theta(u-1/2))+
                    \sinc(\theta(u+1/2))\}^2.
\end{equation}

\begin{lemma}\label{lem:majorant}
The nonnegative function $g$ satisfies $g\ge f$ on $\R$,
$\widehat g(t)=0$ for $|t|\ge\theta$, and
\begin{equation}\label{eq:mass}
 \Lambda:=\int_\R g(u)\,du
 =\frac{2\gamma}{\theta}(1+\sinc\theta)
 \le\frac{61}{32}<2 .
\end{equation}
\end{lemma}
\begin{proof}
Each shifted sinc has a rectangular Fourier transform supported in
$[-\theta/2,\theta/2]$. The square in \eqref{eq:majorant}
therefore has Fourier support $[-\theta,\theta]$.
The convolution is continuous and zero at both boundary points,
so the conclusion includes equality $|t|=\theta$.
Parseval gives the displayed formula for $\Lambda$.
As $\sin(4\pi/5)=\sin(\pi/5)\le\pi/5$,
$\sinc(4/5)\le1/4$, proving its stated upper bound.

The remaining assertion $g\ge f$ is proved by the small exact
rational certificate in Appendix~\ref{app:certificate}.
Specifically it establishes
\[
 g(k/256)-f(k/256)>9/100 \quad(0\le k\le128),
 \qquad \sup_{u\in I}|(g-f)'(u)|<16.
\]
Both functions are even, and every point of $[0,1/2]$ is at
distance at most $1/512$ from one of these nodes. Hence on $I$
\[
 g(u)-f(u)>9/100-16/512=47/800>0.
\]
Outside $I$ the assertion follows from $f=0$ and $g\ge0$.
The derivative estimate is the essential continuous bridge;
the grid inequalities alone would not suffice.
\end{proof}

\begin{lemma}[Close-pair reward]\label{lem:close}
For $|t|\le4/5$, one has $K(t)>134/825$ and
$K(t)^2>1/40>c$.
\end{lemma}
\begin{proof}
For $0\le t\le4/5$, the functions $\cos(\sqrt2u)$ and
$\cos(2\pi tu)$ both decrease on $[0,1/2]$.
Chebyshev's integral covariance inequality, divided by $Z_0$,
gives for the unperturbed kernel $K_0(t)\ge\sinc t\ge\sinc(4/5)$.
The last monotonicity follows by differentiating $\sin x/x$ on
$[0,\pi]$. Using $\pi^2<10$ and
$\sin x\ge x-x^3/6$ at $x=\pi/5$ gives
$\sinc(4/5)>7/30$.
The perturbation of the normalized kernel has modulus at most
$S/Z_0<39/550$. Thus
$K(t)>7/30-39/550=134/825$.
Finally $(134/825)^2=17956/680625>1/40>c$.
Evenness covers negative $t$.
\end{proof}

\section{Lossless spectral assembly}

Choose even $\chi_\varepsilon\in C_c^\infty((-1/2,1/2))$,
$0\le\chi_\varepsilon\le1$, tending to one on the interior, and set
\[
 m_\varepsilon=\int_I f\chi_\varepsilon^2,\qquad
 f_\varepsilon=f\chi_\varepsilon^2/m_\varepsilon,\qquad
 K_\varepsilon=\widehat f_\varepsilon,\qquad
 d_\varepsilon=\|f_\varepsilon-f\|_1.
\]
Then $m_\varepsilon\to1$, $d_\varepsilon\to0$, and
$f_\varepsilon\to f$ in $L^1\cap L^2$.
The square root $\eta_\varepsilon=\sqrt{f_\varepsilon}$ is
smooth, even, compactly supported in the interior of $I$.
On the real axis, $|K_\varepsilon-K|\le d_\varepsilon$
and both kernels have modulus at most one.
Throughout the argument fix $\varepsilon$ sufficiently small that
$m_\varepsilon>61/64$.

For a finite set $Y$ of $n$ real points let
$U=(K_\varepsilon(y_i-y_j))_{i,j}$ be its positive Gram matrix,
with unit diagonal. Define
\[
 \phi_2(t)=t^2-(t-2)_+^2,\qquad
 J(U)=\tr\phi_2(U)-n .
\]
The scalar function $\phi_2$ is convex on $[0,\infty)$.

\begin{proposition}\label{prop:lossless}
For every finite $Y$,
\begin{equation}\label{eq:gain}
 J(U)\ge cn-B\,\operatorname{span}(Y)-7c-42d_\varepsilon n.
\end{equation}
\end{proposition}
\begin{proof}
Select a maximal collection of disjoint pairs with distance
strictly less than $\theta$. The remaining set $S$ is
$\theta$-separated. For any coefficient vector $(z_j)$ on $S$,
the pointwise majorant gives
\[
 \int f_\varepsilon(u)\left|\sum_{j\in S}z_je^{-2\pi iy_ju}\right|^2du
 \le\frac1{m_\varepsilon}
 \int g(u)\left|\sum_{j\in S}z_je^{-2\pi iy_ju}\right|^2du
 =\frac{\Lambda}{m_\varepsilon}\sum_{j\in S}|z_j|^2 .
\]
The cross terms vanish even at separation exactly $\theta$.
Thus $0\preceq U_{SS}\preceq(\Lambda/m_\varepsilon)I\prec2I$.
Normalization by $m_\varepsilon$ is necessary here.
Consequently $\phi_2(U_{SS})=U_{SS}^2$.

Write $m=|S|$. Apply each local certificate to $K$ and transfer
only its finitely many weighted pair terms to $K_\varepsilon$.
Each square changes by at most $3d_\varepsilon$; the total
weight in one window is at most $14$ by \eqref{eq:capacity}.
The sliding argument, including the cases $m\le7$, yields
\[
 \tr\phi_2(U_{SS})-m
 \ge cm-B\,\operatorname{span}(S)-7c-42d_\varepsilon m .
\]
Each removed pair has a Gram matrix with eigenvalues
$1\pm|K_\varepsilon(t)|$ in $[0,2]$.
Its contribution to $J$ is
$2K_\varepsilon(t)^2\ge2c-6d_\varepsilon$, by
Lemma~\ref{lem:close}.

Diagonalize the principal block for $S$ and each two-point block.
Scalar Jensen in the resulting orthonormal basis proves
$\tr\phi_2(U)\ge\sum_{\text{blocks}}\tr\phi_2(U_{\text{block}})$.
This requires scalar convexity, not operator convexity.
Summing the pair rewards and the separated-set reward gives
\eqref{eq:gain}, since $\operatorname{span}(S)\le
\operatorname{span}(Y)$ and the pair error is at most
$3d_\varepsilon(n-m)\le42d_\varepsilon(n-m)$.
Empty or small separated sets are covered by the same single
boundary charge $7c$.
\end{proof}

\section{All-zero Hermitian counting and the energy constant}

\subsection{The unconditional analytic input}

For each fixed $\varepsilon$, the unconditional correlation formula
of Baluyot, Goldston, Suriajaya and Turnage-Butterbaugh
\cite[Lemma~5]{BGST}, with the correction described in
\cite[Section~3]{BGSTcorr}, gives
\begin{equation}\label{eq:unweighted}
 \sum_{0<\Im\rho,\Im\rho'\le T}
 K_\varepsilon\left(i(\rho-\rho')\frac{\log T}{2\pi}\right)^2
 =(C(f_\varepsilon)+o_\varepsilon(1))N(T),
\end{equation}
where
\begin{equation}\label{eq:energy}
 C(h)=\int_Ih(u)^2du+
           \int_I\int_I|u-v|h(u)h(v)\,du\,dv.
\end{equation}
The sum includes multiplicities and all off-line zeros.
For completeness, the weighted lemma applied to a fixed real even
test $q$ supported in $[-1,1]$ has weight
$w(\rho-\rho')=4/(4-(\rho-\rho')^2)$ and main term
$q(0)+2\int_0^1u q(u)\,du$.
Apply it separately to the two fixed smooth functions
$q=f_\varepsilon*f_\varepsilon$ and $q''$.
For $z=i(\rho-\rho')\log T/(2\pi)$, integration by parts gives
\[
 \left(\widehat q(z)-\frac{\widehat{q''}(z)}
                            {4(\log T)^2}\right)w(\rho-\rho')
 =K_\varepsilon(z)^2.
\]
The second fixed-test contribution is $O_\varepsilon(N/(\log T)^2)$;
the first gives \eqref{eq:unweighted}. The lemma allows signed tests,
so its use for $q''$ is legitimate.
This is the fixed-smoothing argument also used in \cite{Lamzouri}.
No estimate uniform in a $T$-dependent smoothing parameter is used.

\subsection{Retaining off-line zeros}

Put $L=\log T$ and $z_\rho=-i(\rho-1/2)L/(2\pi)$.
The multiset of these points is conjugation-invariant.
First in the complex Hilbert space $L^2(I;\mathbb C)$ set
$v_z(u)=\eta_\varepsilon(u)e^{-2\pi izu}$.
For each distinct nonreal conjugate pair write
$g_z=(v_z+v_{\bar z})/2$ and
$h_z=(v_z-v_{\bar z})/(2i)$.
The vectors $v_x$ for real $x$, together with $g_z$ and $h_z$,
belong to the real Hilbert space
$\mathcal H_{\R}=\{v:\overline{v(u)}=v(-u)\}$:
the conjugate reflection $J$ satisfies $Jv_z=v_{\bar z}$.
Define the finite self-adjoint operator on their real span by
\[
 A_T=\sum_{\text{distinct real }x}\mu_xv_x\otimes v_x+
       2\sum_{\text{distinct nonreal pairs}}\mu_z
                  (g_z\otimes g_z-h_z\otimes h_z).
\]
It acts on the finite span of these real-type vectors.
The identity $\|g_z\|^2-\|h_z\|^2=1$ and direct tensor expansion give
\begin{equation}\label{eq:operator}
 \tr A_T=N(T),\qquad
 \tr A_T^2=(C(f_\varepsilon)+o_\varepsilon(1))N(T).
\end{equation}
Indeed the Hilbert--Schmidt square expands to
$\sum_{z,w}K_\varepsilon(z-\bar w)^2$; conjugation reindexes this
as the sum in \eqref{eq:unweighted}. Individual complex terms
need not be positive.

Let $P$ be the sum of the unit features at simple real points,
$n=N_0^s(T)$, and $Q=A_T-P$. If $r$ is the number of repeated
real points and $k$ the number of nonreal conjugate pairs, then
\[
 n_+(Q)\le b:=r+k,\qquad N(T)\ge n+2b.
\]
Each remaining real point has multiplicity at least two; each
nonreal pair contributes at least two zeros and at most one
positive rank direction. No location-based deletion has occurred.

\begin{lemma}[Finite spectral counting]\label{lem:counting}
If $A=P+Q$, $P$ is the sum of $n$ unit features with Gram $U$,
$n_+(Q)\le b$, $\tr A=N$ and $N\ge n+2b$, then
\begin{equation}\label{eq:counting}
 \tr A^2\ge2(N-n)+\tr\phi_2(U)
           =2N-n+J(U).
\end{equation}
\end{lemma}
\begin{proof}
Enlarge the finite ambient space by zero directions if necessary.
If $\lambda_i(P)=p_i$ are its $n$ Gram eigenvalues, including
zeros, minmax gives $\lambda_{b+i}(A)\le p_i$.
For $t\le p$, $p\ge0$,
$4t-t^2\le4p-\phi_2(p)$; also $4t-t^2\le4$ for all real $t$.
After the first $b$ eigenvalues use the first inequality; any
remaining eigenvalues are nonpositive and contribute at most zero.
Since $\sum p_i=n$, this proves
$4N-\tr A^2\le4b+4n-\tr\phi_2(U)$.
Now $4b\le2(N-n)$ gives \eqref{eq:counting}.
Zeros in the $n$-dimensional Gram spectrum are included once;
the added ambient zero directions contribute nothing.
\end{proof}

\subsection{Energy and completion of the proof}

Define $(\mathcal Th)(u)=h(u)+\int_I|u-v|h(v)\,dv$.
For $f_0=\cos(\sqrt2u)/Z_0$, $(\mathcal Tf_0)''=0$ and evenness
makes $\mathcal Tf_0$ constant. Therefore a zero-mass perturbation
has zero cross energy with $f_0$. Orthogonality of the integer
cosines, and
$\mathcal T\cos(2\pi ju)=(1-1/(2\pi^2j^2))\cos(2\pi ju)
+\text{constant}$, give
\begin{equation}\label{eq:amenergy}
 C(f)=\frac12+\beta\cot\beta+
 \frac{\displaystyle\sum_{j=1}^{12}c_j^2
           (1/2-1/(4\pi^2j^2))}{2\sin^2\beta}.
\end{equation}
The exact rational computation in Appendix~\ref{app:certificate}
establishes
\begin{equation}\label{eq:gainconstant}
 2-C(f)>\frac{67216841}{10^8}.
\end{equation}
In particular the energy cost of the perturbation is paid;
the unperturbed energy constant is not substituted.

\begin{proof}[Proof of Theorem~\ref{thm:main}]
The simple real positions lie in $(0,X_T]$, with
$X_T=T\log T/(2\pi)$ and $X_T/N(T)\to1$ by
Riemann--von Mangoldt. Apply Proposition~\ref{prop:lossless}
to their Gram and then Lemma~\ref{lem:counting}:
\[
 (1-c+42d_\varepsilon)n
 \ge2N(T)-\tr A_T^2-BX_T-7c.
\]
Fix sufficiently small $\varepsilon$, let $T\to\infty$ using
\eqref{eq:operator}, and only then let $\varepsilon\to0$.
The $L^1\cap L^2$ convergence implies
$C(f_\varepsilon)\to C(f)$ in \eqref{eq:energy}.
Consequently
\[
 \liminf_T\frac{n}{N(T)}
 \ge\frac{2-C(f)-B}{1-c}
 >\frac{67216841-404350}{10^8-805003}
 =\frac{66812491}{99194997}.
\]
The weak inequality stated in the theorem is therefore valid.
The boundary charge $7c$ disappears only after division by $N(T)$.
\end{proof}

\section{Verification scope and further improvement}\label{sec:trust}

Input~\ref{in:pc8} is pinned to the public
\texttt{dani-bound-2026} source \cite{PC8} at commit
\[
\texttt{d272437e7ebeb17b87c8c3d7cece592aa23785ab}.
\]
Its raw \texttt{Solution.lean} is 1,675,641 bytes with SHA-256
\[
\texttt{\footnotesize 012c6ac5f9282158a686500dc0c967bf0a9b00e1a6192734d8f237bb23dd2d5f}.
\]
The repository admission combines a continuous interval-soundness
audit with all 240 exact integer compiled evaluations, seven complete
finite structural checks, and Lean kernel-checked generic structural
and minorant bridges. It includes the value and derivative tables,
point lists, convexity chains, noncompact gap coverage and all 32
root trees. The evidence is recorded in repository note~483
and its linked audits \cite{Repo}.
The compiled evaluations are not 240 kernel-reduced proofs, and
the full upstream zeta formalization was not rebuilt.
The present result has that explicit finite-computation trust scope.

The new majorant and energy checker uses only Python's standard-library
arbitrary-precision integers and rational arithmetic, with assertions
enabled. It does not run the much larger native search supporting
the numerical headline in \cite{Knaus}. In particular that search's
\texttt{native\_decide} computation axiom is not an input to this paper.
Neither this checker nor the earlier symbolic bridges constitute
a complete Lean verification of Theorem~\ref{thm:main}.
The mathematical transfer, finite computation and formalized
sublemmas are separate evidence.
The manuscript was prepared with AI assistance and internal
independent reviews recorded in the repository; these reviews are
distinct from external peer review.

The mechanism suggests a concrete next problem. Keep the same
window and majorant, but seek a separated seven-gap certificate
with reward $x$, mean gap fee $7\alpha$ and a bounded telescoping
storage correction. Its actual objective is
$(2-C(f)-7\alpha)/(1-x)$, not reward $x$ alone.
All storage endpoint terms must be retained before taking the
limit. No such improved certificate is established here.
The present proportion result also does not improve the project's
separate zero-free-region bounds.

\appendix
\section{Small exact rational certificate}\label{app:certificate}

The accompanying standard-library program is
{\small\path{scripts/am_lossless_majorant_certificate.py}}, with report
{\small\path{output/am-lossless-majorant-certificate.json}}.
From the repository root run
\[
 \texttt{python -B scripts/am\_lossless\_majorant\_certificate.py --check}.
\]
The checker refuses optimized Python execution.
The following specifies the full continuous bridge independently
of floating-point output.

For a nonnegative rational interval $q\subset[0,10)$, it encloses
$\cos\sqrt q$ and $\sin\sqrt q/\sqrt q$ by interval evaluation of
\[
 \sum_{k=0}^{12}\frac{(-1)^kq^k}{(2k)!},
 \qquad
 \sum_{k=0}^{12}\frac{(-1)^kq^k}{(2k+1)!},
\]
with symmetric remainder bounds
$q_{\max}^{13}/26!$ and $q_{\max}^{13}/27!$, respectively.
The alternating tails decrease from the first omitted term.
For the integer cosine modes the rational phase $2ju$ is
reduced modulo two to $[-1,1]$ before multiplication by $\pi$.
No division by a small trigonometric argument is required.
The normalization $Z_0=\sin\beta/\beta$ has the exact argument
squared $q=1/2$.

Machin's identity
$\pi=16\arctan(1/5)-4\arctan(1/239)$, with the first 24 terms
of each alternating arctangent series and its next-term remainder,
certifies
\[
 3141592653/10^9<\pi<3141592654/10^9.
\]
These enclosures, at all 129 nodes $u=k/256$, certify the strict
lower margin $g(u)-f(u)>9/100$.
On all of $I$, $|\sinc'|\le\pi/2$ follows from its integral
representation. Thus
\[
 |(g-f)'|\le4\gamma\pi\theta+
       \frac{\sqrt2+2\pi\sum_jj|c_j|}{Z_0}
 <4\gamma\frac{22}{7}\theta+
       \frac{3/2+(44/7)\sum_jj|c_j|}{9/10}<16.
\]
This proves the global margin $47/800$ stated above.

For \eqref{eq:gainconstant}, use the same Taylor intervals
to bound $\cos\beta/Z_0$ from above, $Z_0^2=2\sin^2\beta$
from below, and $\pi$ from above in the positive correction
in \eqref{eq:amenergy}. All resulting comparisons are exact rational
inequalities. A shorter independent certificate uses
\[
\begin{aligned}
 u&=\frac{3329448031}{4379443200}>\cos\beta,&
 z&=\frac{153778383027779}{167382319104000}<Z_0,\\
 d&=\frac{180493}{213840}<2\sin^2\beta,&
 p&=\frac{5277328977275528}{1679825970703125}>\pi.
\end{aligned}
\]
These arise, respectively, from cosine terms through index 6,
sinc terms through index 7, cosine at $\sqrt2$ through index 6,
and Machin arctangent upper/lower sums through indices 4 and 1.
Direct rational arithmetic gives
\[
 \frac32-\frac uz-
 \frac{\sum_jc_j^2(1/2-1/(4p^2j^2))}{d}
 >\frac{67216841}{10^8}.
\]
The positive margin is greater than $5.6\cdot10^{-10}$.
The checker also verifies every span capacity, the close-pair
constants, the majorant mass bound and
\eqref{eq:comparison}. These small checks do not substitute for
the full seven-dimensional certificate in Input~\ref{in:pc8}.

\begin{thebibliography}{9}
\bibitem{BGST}
S. A. C. Baluyot, D. A. Goldston, A. I. Suriajaya and
C. L. Turnage-Butterbaugh,
\emph{An unconditional Montgomery theorem for pair correlation of
zeros of the Riemann zeta function},
Acta Arith. 214 (2024), 357--376;
\href{https://arxiv.org/abs/2306.04799}{arXiv:2306.04799}.

\bibitem{BGSTcorr}
S. A. C. Baluyot, D. A. Goldston, A. I. Suriajaya and
C. L. Turnage-Butterbaugh,
\emph{Pair correlation of zeros of the Riemann zeta function I:
proportions of simple zeros and critical zeros},
\href{https://arxiv.org/abs/2501.14545v3}{arXiv:2501.14545v3},
Section~3 and its correction of the earlier Theorem~1.

\bibitem{Lamzouri}
Y. Lamzouri,
\emph{A new proof that more than $2/3$ of the zeros of the Riemann
zeta function are simple and on the critical line},
\href{https://arxiv.org/abs/2609.02882v1}{arXiv:2609.02882v1},
Section~3, fixed-smoothing removal of the correlation weight.

\bibitem{PC8}
\emph{AM eight-point PC8CL certificate}, public
\texttt{dani-bound-2026} submission, fixed commit
\texttt{d272437e7ebeb17b87c8c3d7cece592aa23785ab},
\href{https://github.com/josusanmartin/riemann/tree/d272437e7ebeb17b87c8c3d7cece592aa23785ab/submissions/dani-bound-2026/proof}{public proof source}.
The cosine-perturbed window and local weights are prior inputs.

\bibitem{Knaus}
K. M. Knausg{\aa}rd,
\emph{More than 83.9\% of the zeros of the Riemann zeta function
are distinct and more than 67.35\% are simple and on the critical line},
\href{https://arxiv.org/abs/2610.08965v1}{arXiv:2610.08965v1},
October 6, 2026, especially Section~5.

\bibitem{Repo}
RH--Weil research project,
\href{https://github.com/lixiang90/RH-Weil}{public repository},
note~483, the PC8 finite replay and independent continuous audits;
note~498 and the present majorant certificate.
\end{thebibliography}
\end{document}
